\documentclass[12pt]{exam}
\usepackage[margin=0.1in]{geometry}

%\documentclass[12pt]{exam}
\usepackage[utf8]{inputenc}
\usepackage{amsmath,amssymb}
\usepackage{graphicx}
\usepackage{enumitem}
\usepackage{multicol}
\usepackage{xcolor}
\usepackage{tikz}
\usetikzlibrary{arrows.meta,positioning,shapes.callouts}

\newcommand{\tfq}[3]{%
  \question[#1]\textbf{T/F:} #2%
  % Only show the choices when printing answers/keys
  \ifprintanswers
    \par\begin{oneparchoices}
      #3
    \end{oneparchoices}
  \fi
}

\makeatletter
\newcommand{\TFQ@choices}{} % temp holder

\newenvironment{tfqverb}[2]{%
  \def\TFQ@choices{#2}% store choices for use at end
  \question[#1]\textbf{T/F:}\ %
}{%
  \ifprintanswers
    \par\begin{oneparchoices}\TFQ@choices\end{oneparchoices}
  \fi
}
\makeatother

\printanswers
% optional styling for the correct choice:
\CorrectChoiceEmphasis{\bfseries}    % make the correct choice bold
% or e.g. \CorrectChoiceEmphasis{\color{red}\bfseries}

\begin{document}


\begin{center}
\textbf{\Large Review 1 — Database Management Systems}

\textbf{CSC 3300 \quad Spring 2026}
\end{center}

\begin{questions}
%---------------------------------TRUE/FALSE---------------------------------
\section*{Part I — True/False}

{
  \renewcommand{\choicelabel}{}

  \tfq{2}{In a relation instance, all tuples must have non-\emph{null} values for attributes that are part of the primary key.}{\CorrectChoice True \choice False}

  \tfq{2}{The \textbf{takes} relation has exactly five foreign-key attributes.}{\choice True \CorrectChoice False}

  \tfq{2}{In the University Database, a course can have itself as a prerequisite.}{\CorrectChoice True \choice False}

  \tfq{2}{In the \textbf{section} relation, the attributes \(course\_id\) and \(sec\_id\) must always have distinct values from one another.}{\choice True \CorrectChoice False}

  \tfq{2}{In the University Database, the \(course\_id\) can be \texttt{'GoodLuck'}.}{\CorrectChoice True \choice False}

  % ---- USE THE VERBATIM-SAFE ENV ON ITS OWN (not inside \tfq) ----
  \begin{tfqverb}{2}{\CorrectChoice True \choice False}
  When invoked subsequently, the following sequence of commands will execute
  successfully on a database instance with no tuples in any of the tables:
\begin{verbatim}
insert into course  values ('3300', 'DBMS', null, 3);
insert into section values ('3300', '001', 'Fall', 2022, null, null, 'A');
\end{verbatim}
  \end{tfqverb}

\tfq{2}{%
In any instance of the University Database, if the \textbf{instructor} relation contains 0 tuples,
then the \textbf{teaches} relation also contains 0 tuples.
}{\CorrectChoice True \choice False}

\tfq{3}{%
Consider the following modifications to the University Database schema:
\begin{itemize}
  \item Remove the \textbf{teaches} table
  \item Add an attribute \(\mathit{i\_id}\) to the \textbf{section} table
  \item Make \(\mathit{i\_id}\) a foreign key referencing the \textbf{instructor} table
\end{itemize}
\noindent In the modified schema, each section can be taught by at most one instructor.
}{\CorrectChoice True \choice False}

}


%---------------------------------MC---------------------------------
\section*{Part II — Multiple choice (choose one)}
\question[2] Assume that there are no tuples in any of the relation instances of the University Database. 
What is the minimal number of tuples that need to be added to the database before one tuple can be added 
to the table \textbf{teaches}?
\begin{choices}
  \CorrectChoice 3
  \choice 2
  \choice 5
  \choice 4
\end{choices}

\question[3] What is the key structure of the \textbf{section} relation?
\begin{choices}
  \CorrectChoice It has one primary key and two foreign keys.
  \choice It has four primary keys and three foreign keys.
  \choice It has four primary keys and two foreign keys.
  \choice It has one primary key (with four attributes) and three foreign keys.
\end{choices}

%---------------------------------MS---------------------------------
\section*{Part III — Multiple select (select all that apply)}
\question[4] Under the condition that the insertion does not violate any primary key constraints, into which table(s) can a tuple \emph{always} be inserted, regardless of the current database instance?
\begin{checkboxes}
  \CorrectChoice \textbf{course}
  \choice \textbf{advisor}
  \CorrectChoice \textbf{classroom}
  \choice \textbf{prereq}
\end{checkboxes}


\question[4] Which SQL queries retrieve the IDs and titles of \textbf{Comp. Sci.} courses offered in \textbf{Spring 2010}?
\begin{checkboxes}

  \CorrectChoice
\begin{verbatim}
select course.course_id, title
from course, section
where course.course_id = section.course_id
  and course.dept_name = 'Comp. Sci.'
  and section.semester = 'Spring'
  and section.year = 2010;
\end{verbatim}

  \choice
\begin{verbatim}
select course.course_id, title
from course, section
where course.dept_name = 'Comp. Sci.'
  and section.semester = 'Spring'
  and section.year = 2010;
\end{verbatim}

  \choice
\begin{verbatim}
select course.course_id, title
from course, section, department
where course.course_id = section.course_id
  and department.dept_name = 'Comp. Sci.'
  and section.semester = 'Spring'
  and section.year = 2010;
\end{verbatim}

  \choice
\begin{verbatim}
select course_id, title
from course
where dept_name = 'Comp. Sci.'
  and semester = 'Spring'
  and year = 2010;
\end{verbatim}

\end{checkboxes}



\question[4] Which SQL queries retrieve the IDs and names of students along with their advisors' IDs and names?
\begin{checkboxes}

  \choice
\begin{verbatim}
select s.ID, s.name, i.ID, i.name
from student as s, instructor as i, teaches as tc, takes as tk
where tc.course_id = tk.course_id
  and tc.sec_id = tk.sec_id
  and tc.semester = tk.semester
  and tc.year = tk.year
  and tc.ID = i.ID
  and tk.ID = s.ID;
\end{verbatim}

  \CorrectChoice
\begin{verbatim}
select s.ID, s.name, i.ID, i.name
from student as s, instructor as i, advisor as a
where a.s_ID = s.ID
  and a.i_ID = i.ID
  and a.i_ID is not null;
\end{verbatim}

  \CorrectChoice
\begin{verbatim}
select s.ID, s.name, i.ID, i.name
from student as s, instructor as i, advisor as a
where a.s_ID = s.ID
  and a.i_ID = i.ID;
\end{verbatim}

  \choice
\begin{verbatim}
select s.ID, s.name, i.ID, i.name
from student as s, instructor as i, advisor as a, teaches as tc, takes as tk
where a.s_ID = s.ID
  and a.i_ID = i.ID
  and tc.course_id = tk.course_id
  and tc.sec_id = tk.sec_id
  and tc.semester = tk.semester
  and tc.year = tk.year
  and tc.ID = i.ID
  and tk.ID = s.ID;
\end{verbatim}

\end{checkboxes}


\question[4] For the schema below, which statements are true?
\begin{verbatim}
create table instructor (
  id         char(5),
  name       varchar(20),
  dept_name  varchar(20),
  salary     numeric(8,2)
);
\end{verbatim}
\begin{checkboxes}
  \CorrectChoice Creates a table named \texttt{instructor} in the current database (and would fail if a table with that name already exists).
  \choice The attribute \texttt{ID} uniquely identifies a tuple in \texttt{instructor}.
  \CorrectChoice The \texttt{ID} of an instructor can be \texttt{'Hola'}.
  \CorrectChoice The \texttt{name} of an instructor can be \texttt{'Anna M. Brown'}.
\end{checkboxes}

\end{questions}


\end{document}
